%% Section 8: Conclusion

\section{Conclusion}
\label{sec:conclusion}

This paper studies the real effects of Nigeria's 2022--23 naira crunch, an
abrupt, policy-induced reduction in the supply of physical currency, and
the role of digital financial infrastructure in moderating those effects.
The primary evidence comes from a firm-level panel, exploiting pre-shock
cross-state variation in digital payment adoption, measured eleven days before
the \CBN{} announcement, the studys shows large and persistent effects
on enterprise survival and employment. A 45-month pre-shock nighttime-light panel for all 774 Nigerian local
government areas validates the parallel-trends assumption as high- and
low-fintech states followed statistically identical \NTL{} trajectories before
the crunch, and the identifying assumption is robust to pre-trend violations
up to 20 times the observed magnitude.

First, enterprises in higher-fintech states were 1.7 percentage points
more likely to remain in operation at the February 2023 peak crunch, had
0.29 more workers, and the survival advantage widened to 7.7 percentage
points by April 2024. Effects are concentrated among rural and small
enterprises, the units most exclusively dependent on physical cash. 
Pairs bootstrap and randomisation-inference $p$-values confirm the R11 survival
and R7 employment results hold at the 1\% level under all inference methods;
the R7 survival effect is significant at the 10\% level under randomisation
inference and read as supportive evidence of peak-crunch insulation. 
Second, the 2025 \WBES{} cross-section
provides descriptive evidence of elevated cash dependence in low-fintech
zones, consistent with a reinforcing dynamic that the crunch appears to have
entrenched financial exclusion where digital infrastructure was thinnest.

To quantify the aggregate stakes and applying sampling weights to the Round~5
enterprise module yields approximately 20.5~million non-farm household
enterprises in Nigeria at the pre-shock baseline. Closing the fintech gap
between the 10th- and 90th-percentile state (a 1.85 standard-deviation
difference) would have reduced enterprise closures at peak crunch by
approximately 648,000 enterprises and preserved approximately 11~million
enterprise-worker positions during the most acute phase of the shortage.
By April 2024, the cumulative survival advantage translates to approximately
2.9~million fewer enterprise closures relative to a counterfactual of
homogeneous low-fintech adoption.

\paragraph{Policy implications.}
First, the results quantify a
macroeconomic stabilisation benefit of digital payment infrastructure that
is distinct from its better-documented microeconomic benefits for individual
households \citep{jack2014risk,suri2016mobile}. Digital payment systems do
not merely redistribute welfare within an economy; they reduce the aggregate
cost of a cash supply shock to the economy as a whole. This stabilisation
benefit should enter the \CBN's cost--benefit calculus when designing currency
operations, implementing withdrawal limits, or accelerating the rollout of
the e-Naira central bank digital currency.

Second, the results suggest that the sequencing of currency policy matters.
The naira crunch appears to have widened the digital--financial divide, as
firms in low-fintech zones entered a persistent cycle of exclusion. This
implies that currency redesigns or cash-withdrawal restrictions deployed
before digital payment infrastructure has reached a critical mass can have
lasting adverse effects on financial inclusion, the opposite of their stated
objective. Policy makers should ensure adequate digital payment coverage
before constraining cash supply.

\paragraph{Generalisability.}
While the estimates are identified from a single episode in Nigeria, the
mechanism, digital payments as a substitute transaction medium when cash is
scarce, is general. The satellite--survey methodology employed here
is directly transferable to other African economies or low-income settings, 
offering a way for empirical assessment in environments where macroeconomic data are limited, unavailable or lagged.

\paragraph{Limitations.}
First, the \NLPS{} enterprise panel tracks
household enterprises, which are likely more cash-intensive than the
formal-sector firms in the \WBES{}; the two samples are not directly
comparable. Second, the causal interpretation of the long-run \WBES{}
comparison requires the assumption that zone composition did not change
differentially across zones between 2014 and 2025, an assumption this study was unable to
verify from the public-release data.
